% ============================================================================
% PART 07 - CROSS-UNIT RELATIONSHIPS
% ============================================================================

\labelsec{cross_unit_relationships}

The three SHT units were designed to be geometrically orthogonal. This chapter tests that assumption and characterizes the one unexpected exception: a strong coupling between Vij and Mamton.

\section{Pairwise Correlations}

All analyses use the merged dataset of $N = 65$ structures for which all three units are measured. Statistical tests: Pearson correlation with two-tailed t-test; 95\% CI via Fisher z-transformation.

\subsection{Vij vs Mamton: A Fundamental Discovery}

\begin{keyresult}
\begin{equation}
    r(v_0,\,\Omega) = 0.8285, \quad p = 1.633 \times 10^{-17}, \quad 95\%\,\text{CI} = [0.733,\;0.892]
\end{equation}
\begin{equation}
    R^2 = 0.686, \quad t = 11.744, \quad \mathrm{df} = 63
\end{equation}
This is the strongest inter-unit relationship observed in the SHT framework.
\end{keyresult}

The linear fit is:
\begin{equation}\label{eq:vij_mamton_linear}
    \Omega = 86{,}683.4\,v_0 - 8{,}109.9
\end{equation}
The fitted slope predicts that each unit increase in curvature (L$^{-1}$) raises the skewness magnitude by $\sim\!86{,}683$ L$^3$.

\begin{figure}[h]
    \centering
    \includegraphics[width=0.85\textwidth]{pairwise_scatter}
    \caption{Pairwise scatter plots for all three SHT units ($N=65$). The Vij--Mamton panel (top-right) shows the strong linear coupling ($r=0.829$). The other two panels confirm statistical independence.}\labelfig{pairwise_scatter}
\end{figure}

\textbf{Heteroscedasticity note:} Mamton has high variance (std = 2,622, median = 79.7). The log-Mamton vs Vij correlation is also strong: $r(\log\Omega, v_0) = 0.77$, confirming the relationship is not driven solely by the two extreme nucleosome outliers (1AOI, 1F66). Removing both nucleosomes ($N=63$) still yields $r = 0.72$, $p < 10^{-10}$.

\subsection{Vij vs Shambhian: Confirmed Independence}

\begin{keyresult}
    $r(v_0,\,\eta_r) = -0.0851$, \quad $p = 0.500$ \quad ($N=65$)
\end{keyresult}

Not significant. Curvature does not predict cross-sectional eccentricity, confirming that these two units capture different aspects of DNA geometry. A highly curved sequence can be either circular ($\eta_r \approx 1$, like 1I8M) or markedly eccentric ($\eta_r \approx 0.53$, like 1LO1).

\subsection{Shambhian vs Mamton: Weakly Independent}

\begin{keyresult}
    $r(\eta_r,\,\Omega) = -0.1522$, \quad $p = 0.226$ \quad ($N=65$)
\end{keyresult}

Numerical trend: more distorted cross-sections (lower $\eta_r$) weakly associate with lower skewness, but the effect is not significant. Shambhian and Mamton are essentially independent.

\section{Summary: Correlation Matrix}

\begin{table}[h]
    \centering
    \caption{Pairwise Pearson Correlations ($N=65$)}\labeltab{correlation_matrix}
    \begin{tabular}{lccc}
        \toprule
         & \textbf{Vij} & \textbf{Shambhian} & \textbf{Mamton} \\
        \midrule
        \textbf{Vij}       & 1.000   & $-0.085$               & $\mathbf{0.829}^{***}$  \\
        \textbf{Shambhian} & $-0.085$ & 1.000                  & $-0.152$               \\
        \textbf{Mamton}    & $0.829^{***}$ & $-0.152$          & 1.000                  \\
        \bottomrule
        \multicolumn{4}{l}{\small $^{***}$: $p < 10^{-15}$; other pairs: $p > 0.2$ (not significant)}
    \end{tabular}
\end{table}

\begin{figure}[h]
    \centering
    \includegraphics[width=0.55\textwidth]{correlation_matrix}
    \caption{Pearson correlation heatmap for the three SHT units ($N=65$). Only the Vij--Mamton pair is significant ($r=0.829$, $p<10^{-15}$).}\labelfig{correlation_matrix}
\end{figure}

The matrix reveals a $2 \times 2$ block structure: Vij and Mamton form one coupled ``Integrity Mode'', while Shambhian stands alone as a ``Shape Mode''. This is confirmed by PCA.

\section{Principal Component Analysis}

PCA on the three-unit vectors $[v_0, \eta_r, \Omega]^T$ (z-scored per unit, $N=65$) reveals:

\begin{table}[h]
    \centering
    \caption{PCA on Three Geometric Units ($N=65$)}\labeltab{pca_units}
    \begin{tabular}{lll}
        \toprule
        \textbf{Component} & \textbf{Variance Explained} & \textbf{Interpretation} \\
        \midrule
        PC1 & 55.8\% & Integrity Mode: high-loading Vij + Mamton \\
        PC2 & 27.6\% & Shape Mode: high-loading Shambhian \\
        PC3 & 16.6\% & Residual / orthogonal noise \\
        \midrule
        PC1+PC2 & \textbf{83.4\%} & 2D manifold in 3D measurement space \\
        \bottomrule
    \end{tabular}
\end{table}

\begin{figure}[h]
    \centering
    \includegraphics[width=0.75\textwidth]{pca_manifold}
    \caption{PCA manifold of the three SHT units. PC1 (Integrity Mode) captures Vij and Mamton; PC2 (Shape Mode) captures Shambhian. The two PCs together explain 83.4\% of variance, showing DNA geometry lives on a 2D reduced manifold.}\labelfig{pca_manifold}
\end{figure}

The 83.4\% in two components indicates the effective dimensionality of DNA geometry as captured by SHT is two, not three. The third dimension (PC3) adds $~17\%$ independent information, which may encode fine-grained local deformations not captured by the two modes.

\section{Mechanistic Hypothesis: Bending Causes Roughness}

The coupling $r = 0.83$ requires a physical explanation. We propose:

\begin{enumerate}
    \item \textbf{Bending induces differential strain}: When DNA bends (high $v_0$), the outer strand is stretched and the inner strand compressed. This creates strand-to-strand asymmetry.
    \item \textbf{Asymmetric strain $\rightarrow$ backbone asymmetry}: Unequal extension of the two strands displaces backbone atoms asymmetrically relative to the helix axis.
    \item \textbf{Backbone asymmetry $\rightarrow$ radial-moment skewness}: The weighted skewness $S_x = \sum w_i x_i r_i^2$ accumulates when atoms preferentially lie on one side of the helix --- exactly what asymmetric bending produces.
\end{enumerate}

\textbf{Supporting evidence:}
\begin{itemize}
    \item A-tracts: simultaneously low Vij ($0.080$) \textit{and} low Mamton ($70.5$) --- \textit{consistent with mechanism}
    \item Nucleosomes: simultaneously high Vij ($0.304$) \textit{and} high Mamton ($\sim 26{,}366$) --- \textit{extreme end of mechanism}
    \item Cisplatin-adduct 1AHD: high eccentricity ($\eta_r = 0.71$, $z = -4.18\sigma$) but intermediate Vij/Mamton --- \textit{consistent with independent shape distortion not coupled to bending}
\end{itemize}

This mechanistic picture predicts that any agent inducing DNA bending (protein binding, drug intercalation, covalent crosslinks) should increase Mamton proportionally to the induced curvature.

\section{Orthogonality: What It Means for Biology}

The overall structure --- one coupled pair + one independent unit --- is advantageous for biological interpretation:

\begin{itemize}
    \item \textbf{Vij alone}: local inverse radius of curvature (where is the DNA bent?)
    \item \textbf{Mamton alone}: skewness of the backbone distribution (how asymmetric is the bending?)
    \item \textbf{Vij + Mamton together}: confirm and quantify bending (the coupling provides internal consistency check)
    \item \textbf{Shambhian alone}: cross-sectional eccentricity (how non-circular is the helix? --- independent of bending)
    \item \textbf{Shambhian decoupled}: identifies structural deformation not driven by gross bending (e.g., groove compression, protein-induced flattening)
\end{itemize}

